\section{Mathematics}
\label{sec:math}

\paragraph{Plain English.}
Each country still wants a baseline pile of food and feed for each grain,
plus (for US maize) an ethanol pile that does not care about wheat or rice
prices. The food-and-feed pile shrinks when that grain is expensive, and
grows a little when a \emph{substitute} grain is expensive. Then each grain
runs the same Gate~0 market it always did: offers, asks, Armington trade,
warehouse. We compute all three demands from this fortnight's three prices,
then all three markets, then the three new prices (Jacobi).

\paragraph{Flex demand.}
For grain \(g\), country \(i\), step \(t\),

\begin{equation}
  d^g_{i,t}
  = C^{g,\mathrm{flex}}_{i,t}
    \prod_h \Bigl(\frac{p^h_t}{p^h_0}\Bigr)^{\eta_{gh}}
  + C^{g,\mathrm{ind}}_{i,t}.
  \label{eq:flex}
\end{equation}

Own exponents \(\eta_{gg}=\varepsilon_g\) from Gate~0
\texttt{CropParams.elast}. Cross exponents (\(g\neq h\))

\begin{equation}
  \eta_{gh}=\sigma\,\rho_{gh}\,|\varepsilon_g|.
  \label{eq:cross}
\end{equation}

At \(p=p_0\) the product is 1 for any \(\sigma\), so the calm twin stays
flat. Code: \(\mathtt{fac}=\exp(\eta\log(p/p_0))\) in
\texttt{sheaf/dynamic\_coupled.py}. Grain order is wheat, rice, maize.

\paragraph{Worked \(\eta\) at \(\sigma=0.6\).}

\[
\eta =
\begin{pmatrix}
-0.15 & 0.027 & 0.036 \\
 0.036 & -0.20 & 0.024 \\
 0.060 & 0.030 & -0.25
\end{pmatrix}
\quad
\text{(rows/cols: wheat, rice, maize).}
\]

Maize answers wheat more than wheat answers maize because
\(|\varepsilon_{\mathrm{maize}}|>|\varepsilon_{\mathrm{wheat}}|\). That
asymmetry is a feature of~\eqref{eq:cross}, not a coding error.

\paragraph{What is not in this layer.}
Joint multi-grain quadratic program; annual integrable \(M_i\); endogenous
export game; endogenous trade network. Those are other papers
\citep{armington1969}.
